% ======================================================================
% Front matter: Preface, Notations and Conventions, Editor's Foreword
% (transcribed from the 1995 Westview Press edition)
% ======================================================================
\chapter*{Preface}
\addcontentsline{toc}{chapter}{Preface}

Quantum field theory is a set of ideas and tools that combines three of the
major themes of modern physics: the quantum theory, the field concept, and
the principle of relativity.  Today, most working physicists need to know some
quantum field theory, and many others are curious about it.  The theory
underlies modern elementary particle physics, and supplies essential tools to
nuclear physics, atomic physics, condensed matter physics, and astrophysics.
In addition, quantum field theory has led to new bridges between physics and
mathematics.

One might think that a subject of such power and widespread application
would be complex and difficult.  In fact, the central concepts and techniques
of quantum field theory are quite simple and intuitive.  This is especially
true of the many pictorial tools (Feynman diagrams, renormalization group
flows, and spaces of symmetry transformations) that are routinely used by
quantum field theorists.  Admittedly, these tools take time to learn, and
tying the subject together with rigorous proofs can become extremely
technical.  Nevertheless, we feel that the basic concepts and tools of
quantum field theory can be made accessible to all physicists, not just an
elite group of experts.

A number of earlier books have succeeded in making parts of this subject
accessible to students.  The best known of these is the two-volume text
written in the 1960s by our Stanford colleagues Bjorken and Drell.  In our
opinion, their text contains an ideal mixture of abstract formalism,
intuitive explanations, and practical calculations, all presented with great
care and clarity.  Since the 1960s, however, the subject of quantum field
theory has developed enormously, both in its conceptual framework (the
renormalization group, new types of symmetries) and in its areas of
application (critical exponents in condensed matter systems, the standard
model of elementary particle physics).  It is long overdue that a textbook
of quantum field theory should appear that provides a complete survey of
the subject, including these newer developments, yet still offers the same
accessibility and depth of treatment as Bjorken and Drell.  We have written
this book with that goal in mind.

\section*{An Outline of the Book}

This textbook is composed of three major sections.  The first is mainly
concerned with the quantum theory of electromagnetism, which provided the
first example of a quantum field theory with direct experimental
applications.  The third part of the book is mainly concerned with the
particular quantum field theories that appear in the standard model of
particle interactions.  The second part of the book is a bridge between
these two subjects; it is intended to introduce some of the very deep
concepts of quantum field theory in a context that is as straightforward
as possible.

Part~I begins with the study of fields with linear equations of motion,
that is, fields without interactions.  Here we explore the combined
implications of quantum mechanics and special relativity, and we learn how
particles arise as the quantized excitations of fields.  We then introduce
interactions among these particles and develop a systematic method of
accounting for their effects.  After this introduction, we carry out
explicit computations in the quantum theory of electromagnetism.  These
illustrate both the special features of the behavior of electrons and
photons and some general aspects of the behavior of interacting quantum
fields.

In several of the calculations in Part~I, naive methods lead to infinite
results.  The appearance of infinities is a well-known feature of quantum
field theory.  At times, it has been offered as evidence for the
inconsistency of quantum field theory (though a similar argument could be
made against the classical electrodynamics of point particles).  For a long
time, it was thought sufficient to organize calculations in such a way that
no infinities appear in quantities that can be compared directly to
experiment.  However, one of the major insights of the more recent
developments is that these formal infinities actually contain important
information that can be used to predict the qualitative behavior of a
system.  In Part~II of the book, we develop this theory of infinities
systematically.  The development makes use of an analogy between
quantum-mechanical and thermal fluctuations, which thus becomes a bridge
between quantum field theory and statistical mechanics.  At the end of
Part~II we discuss applications of quantum field theory to the theory of
phase transitions in condensed matter systems.

Part~III deals with the generalizations of quantum electrodynamics that
have led to successful models of the forces between elementary particles.
To derive these generalizations, we first analyze and generalize the
fundamental symmetry of electrodynamics, then work out the consequences of
quantizing a theory with this generalized symmetry.  This analysis leads
to intricate and quite nontrivial applications of the concepts introduced
earlier.  We conclude Part~III with a presentation of the standard model
of particle physics and a discussion of some of its experimental tests.

The Epilogue to the book discusses qualitatively the frontier areas of
research in quantum field theory and gives references that can guide a
student to the next level of study.

Where a choice of viewpoints is possible, we have generally chosen to
explain ideas in language appropriate to the applications to elementary
particle physics.  This choice reflects our background and research
interests.  It also reflects our strongly held opinion, even in this age
of intellectual relativism, that there is something special about
unraveling the behavior of Nature at the deepest possible level.  We are
proud to take as our subject the structure of the fundamental
interactions, and we hope to convey to the reader the grandeur and
continuing vitality of this pursuit.

\section*{How to Use This Book}

This book is an introduction to quantum field theory.  By this we mean,
first and foremost, that we assume no prior knowledge of the subject on
the part of the reader.  The level of this book should be appropriate for
students taking their first course in quantum field theory, typically
during the second year of graduate school at universities in the United
States.  We assume that the student has completed graduate-level courses
in classical mechanics, classical electrodynamics, and quantum mechanics.
In Part~II we also assume some knowledge of statistical mechanics.  It is
not necessary to have mastered every topic covered in these courses,
however.  Crucially important prerequisites include the Lagrangian and
Hamiltonian formulations of dynamics, the relativistic formulation of
electromagnetism using tensor notation, the quantization of the harmonic
oscillator using ladder operators, and the theory of scattering in
nonrelativistic quantum mechanics.  Mathematical prerequisites include an
understanding of the rotation group as applied to the quantum mechanics of
spin, and some facility with contour integration in the complex plane.

Despite being an ``introduction'', this book is rather lengthy.  To some
extent, this is due to the large number of explicit calculations and
worked examples in the text.  We must admit, however, that the total
number of topics covered is also quite large.  Even students specializing
in elementary particle theory will find that their first research projects
require only a part of this material, together with additional,
specialized topics that must be gleaned from the research literature.
Still, we feel that students who want to become experts in elementary
particle theory, and to fully understand its unified view of the
fundamental interactions, should master every topic in this book.  Students
whose main interest is in other fields of physics, or in particle
experimentation, may opt for a much shorter ``introduction'', omitting
several chapters.

The senior author of this book once did succeed in ``covering'' 90\% of
its content in a one-year lecture course at Stanford University.  But this
was a mistake; at such a pace, there is not enough time for students of
average preparation to absorb the material.  Our saner colleagues have
found it more reasonable to cover about one Part of the book per semester.
Thus, in planning a one-year course in quantum field theory, they have
chosen either to reserve Part~III for study at a more advanced level or to
select about half of the material from Parts~II and~III, leaving the rest
for students to read on their own.

We have designed the book so that it can be followed from cover to cover,
introducing all of the major ideas in the field systematically.
Alternatively, one can follow an accelerated track that emphasizes the
less formal applications to elementary particle physics and is sufficient
to prepare students for experimental or phenomenological research in that
field.  Sections that can be omitted from this accelerated track are
marked with an asterisk in the Table of Contents; none of the unmarked
sections depend on this more advanced material.  Among the unmarked
sections, the order could also be varied somewhat: Chapter~10 does not
depend on Chapters~8 and~9; Section~11.1 is not needed until just before
Chapter~20; and Chapters~20 and~21 are independent of Chapter~17.

Those who wish to study some, but not all, of the more advanced sections
should note the following table of dependencies:
\begin{center}
\begin{tabular}{ll}
\toprule
Before reading \ldots & one should read all of \ldots\\
\midrule
Chapter 13            & Chapters 11, 12\\
Section 16.6          & Chapter 11\\
Chapter 18            & Sections 12.4, 12.5, 16.5\\
Chapter 19            & Sections 9.6, 15.3\\
Section 19.5          & Section 16.6\\
Section 20.3          & Sections 19.1--19.4\\
Section 21.3          & Chapter 11\\
\bottomrule
\end{tabular}
\end{center}

Within each chapter, the sections marked with an asterisk should be read
sequentially, except that Sections~16.5 and~16.6 do not depend on~16.4.

A student whose main interest is in statistical mechanics would want to
read the book sequentially, confronting the deep formal issues of Part~II
but ignoring most of Part~III, which is mainly of significance to
high-energy phenomena.  (However, the material in Chapters~15 and~19, and
in Section~20.1, does have beautiful applications in condensed matter
physics.)

We emphasize to all students the importance of working actively with the
material while studying.  It probably is not possible to understand any
section of this book without carefully working out the intermediate steps
of every derivation.  In addition, the problems at the end of each chapter
illustrate the general ideas and often apply them in nontrivial, realistic
contexts.  However, the most illustrative exercises in quantum field
theory are too long for ordinary homework problems, being closer to the
scale of small research projects.  We have provided one of these lengthy
problems, broken up into segments with hints and guidance, at the end of
each of the three Parts of the book.  The volume of time and paper that
these problems require will be well invested.

At the beginning of each Part we have included a brief ``Invitation''
chapter, which previews some of the upcoming ideas and applications.
Since these chapters are somewhat easier than the rest of the book, we
urge all students to read them.

\section*{What This Book Is Not}

Although we hope that this book will provide a thorough grounding in
quantum field theory, it is in no sense a complete education.  A dedicated
student of physics will want to supplement our treatment in many areas.
We summarize the most important of these here.

First of all, this is a book about theoretical methods, not a review of
observed phenomena.  We do not review the crucial experiments that led to
the standard model of elementary particle physics or discuss in detail
the more recent experiments that have confirmed its predictions.  Similarly,
in the chapters that deal with applications to statistical mechanics, we
do not discuss the beautiful and varied experiments on phase transitions
that led to the confirmation of field theory models.  We strongly encourage
the student to read, in parallel with this text, a modern presentation of
the experimental development of each of these fields.

Although we present the elementary aspects of quantum field theory in full
detail, we state some of the more advanced results without proof.  For
example, it is known rigorously, to all orders in the standard expansion
of quantum electrodynamics, that formal infinities can be removed from all
experimental predictions.  This result, known as \emph{renormalizability},
has important consequences, which we explore in Part~II.  We do not present
the general proof of renormalizability.  However, we do demonstrate
renormalizability explicitly in illustrative, low-order computations, we
discuss intuitively the issues that arise in the complete proof, and we
give references to a more complete demonstration.  More generally, we have
tried to motivate the most important results (usually through explicit
examples) while omitting lengthy, purely technical derivations.

Any introductory survey must classify some topics as beyond its scope.
Our philosophy has been to include what can be learned about quantum field
theory by considering weakly interacting particles and fields, using
series expansions in the strength of the interaction.  It is amazing how
much insight one can obtain in this way.  However, this definition of our
subject leaves out the theory of bound states, and also phenomena
associated with nontrivial solutions to nonlinear field equations.  We give
a more complete listing of such advanced topics in the Epilogue.

Finally, we have not attempted in this book to give an accurate record of
the history of quantum field theory.  Students of physics do need to
understand the history of physics, for a number of reasons.  The most
important is to acquire a precise understanding of the experimental basis
of the subject.  A second important reason is to gain an idea of how
science progresses as a human endeavor, how ideas develop as small steps
taken by individuals to become the major achievements of the community as
a whole.%
\footnote{The history of the development of quantum field theory and
particle physics has recently been reviewed and debated in a series of
conference volumes: \emph{The Birth of Particle Physics}, L.~M.~Brown and
L.~Hoddeson, eds.\ (Cambridge University Press, 1983); \emph{Pions to
Quarks}, L.~M.~Brown, M.~Dresden, and L.~Hoddeson, eds.\ (Cambridge
University Press, 1989); and \emph{The Rise of the Standard Model},
L.~M.~Brown, M.~Dresden, L.~Hoddeson, and M.~Riordan, eds.\ (Cambridge
University Press, 1995).  The early history of quantum electrodynamics is
recounted in a fascinating book by Schweber (1994).}

In this book we have not addressed either of these needs.  Rather, we have
included only the kind of mythological history whose purpose is to motivate
new ideas and assign names to them.  A principle of physics usually has a
name that has been assigned according to the community's consensus on who
deserves credit for its development.  Usually the real credit is only
partial, and the true historical development is quite complex.  But the
clear assignment of names is essential if physicists are to communicate
with one another.

Here is one example.  In Section~17.5 we discuss a set of equations
governing the structure of the proton, which are generally known as the
Altarelli--Parisi equations.  Our derivation uses a method due to Gribov
and Lipatov (GL).  The original results of GL were rederived in a more
abstract language by Christ, Hasslacher, and Mueller (CHM).  After the
discovery of the correct fundamental theory of the strong interactions
(QCD), Georgi, Politzer, Gross, and Wilczek (GPGW) used the technique of
CHM to derive formal equations for the variation of the proton structure.
Parisi gave the first of a number of independent derivations that converted
these equations into a useful form.  The combination of his work with that
of GPGW gives the derivation of the equations that we present in
Section~18.5.  Dokshitzer later obtained these equations more simply by
direct application of the method of GL.  Some time later, but
independently, Altarelli and Parisi obtained these equations again by the
same route.  These last authors also popularized the technique, explaining
it very clearly, encouraging experimentalists to use the equations in
interpreting their data, and prodding theorists to compute the systematic
higher-order corrections to this picture.  In Section~17.5 we have
presented the shortest path to the end of this convoluted historical road
and hung the name ``Altarelli--Parisi'' on the final result.

There is a fourth reason for students to read the history of physics:
Often the original breakthrough papers, though lacking a textbook's
advantages of hindsight, are filled with marvelous personal insights.  We
strongly encourage students to go back to the original literature whenever
possible and see what the creators of the field had in mind.  We have tried
to aid such students with references provided in footnotes.  Though
occasionally we refer to papers merely to give credit, most of the
references are included because we feel the reader should not miss the
special points of view that the authors put forward.

\section*{Acknowledgments}

The writing of this book would not have been possible without the help and
encouragement of our many teachers, colleagues, and friends.  It is our
great pleasure to thank these people.

Michael has been privileged to learn field theory from three of the
subject's contemporary masters---Sidney Coleman, Steven Weinberg, and
Kenneth Wilson.  He is also indebted to many other teachers, including
Sidney Drell, Michael Fisher, Kurt Gottfried, John Kogut, and Howard
Georgi, and to numerous co-workers, in particular, Orlando Alvarez, John
Preskill, and Edward Witten.  His association with the laboratories of
high-energy physics at Cornell and Stanford, and discussions with such
people as Gary Feldman, Martin Perl, and Morris Swartz, have also shaped
his viewpoint.  To these people, and to many other people who have taught
him points of physics over the years, Michael expresses his gratitude.

Dan has learned field theory from Savas Dimopoulos, Leonard Susskind,
Richard Blankenbecler, and many other instructors and colleagues, to whom
he extends thanks.  For his broader education in physics he is indebted to
many teachers, but especially to Thomas Moore.  In addition, he is grateful
to all the teachers and friends who have criticized his writing over the
years.

For their help in the construction and writing of this book, many people
are due more specific thanks.  This book originated as graduate courses
taught at Cornell and at Stanford, and we thank the students in these
courses for their patience, criticism, and suggestions.  During the years
in which this book was written, we were supported by the Stanford Linear
Accelerator Center and the Department of Energy, and Dan also by the Pew
Charitable Trusts and the Departments of Physics at Pomona College,
Grinnell College, and Weber State University.  We are especially grateful
to Sidney Drell for his support and encouragement.  The completion of this
book was greatly assisted by a sabbatical Michael spent at the University
of California, Santa Cruz, and he thanks Michael Nauenberg and the physics
faculty and students there for their gracious hospitality.

We have circulated preliminary drafts of this book widely and have been
continually encouraged by the response to these drafts and the many
suggestions and corrections they have elicited.  Among the people who have
helped us in this way are Curtis Callan, Edward Farhi, Jonathan Feng,
Donald Finnell, Paul Frampton, Jeffrey Goldstone, Brian Hatfield, Barry
Holstein, Gregory Kilcup, Donald Knuth, Barak Kol, Phillip Nelson,
Matthias Neubert, Yossi Nir, Arvind Rajaraman, Pierre Ramond, Katsumi
Tanaka, Xerxes Tata, Arkady Vainshtein, Rudolf von B\"unau, Shimon
Yankielowicz, and Eran Yehudai.  For help with more specific technical
problems or queries, we are grateful to Michael Barnett, Marty
Breidenbach, Don Groom, Toichiro Kinoshita, Paul Langacker, Wu-Ki Tung,
Peter Zerwas, and Jean Zinn-Justin.  We thank our Stanford colleagues
James Bjorken, Stanley Brodsky, Lance Dixon, Bryan Lynn, Helen Quinn,
Leonard Susskind, and Marvin Weinstein, on whom we have tested many of the
explanations given here.  We are especially grateful to Michael Dine,
Martin Einhorn, Howard Haber, and Renata Kallosh, and to their long-
suffering students, who have used drafts of this book as textbooks over a
period of several years and have brought us back important and
illuminating criticism.

We pause to remember two friends of this book who did not see its
completion---Donald Yennie, whose long career made him one of the heroes
of Quantum Electrodynamics, and Brian Warr, whose brilliant promise in
theoretical physics was cut short by AIDS.

We thank our editors at Addison-Wesley---Allan Wylde, Barbara Holland,
Jack Repcheck, Heather Mimnaugh, and Jeffrey Robbins---for their
encouragement of this project, and for not losing faith that it would be
completed.  We are grateful to Mona Zeftel and Lynne Reed for their advice
on the formatting of this book, and to Jeffrey Towson and Cristine
Jennings for their skillful preparation of the illustrations.  The final
typeset pages were produced at the WSU Printing Services Department, with
much help from Allan Davis.  Of course, the production of this book would
have been impossible without the many individuals who have kept our
computers and network links up and running.

We thank the many friends whose support we have relied on during the long
course of this project.  We are particularly grateful to our parents,
Dorothy and Vernon Schroeder and Gerald and Pearl Peskin.  Michael also
thanks Valerie, Thomas, and Laura for their love and understanding.

Finally, we thank you, the reader, for your time and effort spent studying
this book.  Though we have tried to cleanse this text of conceptual and
typographical errors, we apologize in advance for those that have slipped
through.  We will be glad to hear your comments and suggestions for
further improvements in the presentation of quantum field theory.

\vspace{2em}
\noindent Michael E. Peskin\\
Daniel V. Schroeder

% ======================================================================
\chapter*{Notations and Conventions}
\addcontentsline{toc}{chapter}{Notations and Conventions}

\section*{Units}

We will work in ``God-given'' units, where
\begin{equation}
  \hbar = c = 1 .
\end{equation}
The mass $m$ of a particle is therefore equal to its rest energy $mc^2$,
and also to its inverse Compton wavelength $mc/\hbar$.  For example,
\begin{equation}
  m_\mathrm{electron} = 9.109 \times 10^{-28}\,\mathrm{g}
    = 0.511\ \mathrm{MeV}
    = (3.862 \times 10^{-11}\,\mathrm{cm})^{-1} .
\end{equation}
A selection of other useful numbers and conversion factors is given in the
Appendix.

\section*{Relativity and Tensors}

Our conventions for relativity follow Jackson (1975), Bjorken and Drell
(1964, 1965), and nearly all recent field theory texts.  We use the metric
tensor
\begin{equation}
  g_{\mu\nu} = \begin{pmatrix}
    1 & 0 & 0 & 0\\
    0 & -1 & 0 & 0\\
    0 & 0 & -1 & 0\\
    0 & 0 & 0 & -1
  \end{pmatrix},
\end{equation}
with Greek indices running over $0,1,2,3$ or $t,x,y,z$.  Roman indices---
$i,j,$ etc.\---denote only the three spatial components.  Repeated indices
are summed in all cases.  Four-vectors, like ordinary numbers, are denoted
by light italic type; three-vectors are denoted by boldface type; unit
three-vectors are denoted by a light italic label with a hat over it.
For example,
\begin{align}
  x^\mu &= (x^0, \mathbf{x}),\qquad x_\mu = g_{\mu\nu} x^\nu = (x^0, -\mathbf{x});\\
  p\cdot x &= g_{\mu\nu} p^\mu x^\nu = p^0 x^0 - \mathbf{p}\cdot\mathbf{x}.
\end{align}
A massive particle has
\begin{equation}
  p^2 = p_\mu p^\mu = E^2 - |\mathbf{p}|^2 = m^2 .
\end{equation}
Note that the displacement vector $x^\mu$ is ``naturally raised'', while
the derivative operator,
\begin{equation}
  \partial_\mu = \Bigl( \frac{\partial}{\partial t},\ \nabla \Bigr),
\end{equation}
is ``naturally lowered''.

We define the totally antisymmetric tensor $\epsilon^{\mu\nu\rho\sigma}$ so
that
\begin{equation}
  \epsilon^{0123} = +1 .
\end{equation}
Be careful, since this implies $\epsilon_{0123} = -1$ and
$\epsilon^{1230} = -1$.  (This convention agrees with Jackson but not with
Bjorken and Drell.)

\section*{Quantum Mechanics}

We will often work with the Schr\"odinger wavefunctions of single
quantum-mechanical particles.  We represent the energy and momentum
operators acting on such wavefunctions following the usual conventions:
\begin{equation}
  E = i\frac{\partial}{\partial t},\qquad \mathbf{p} = -i\nabla .
\end{equation}
These equations can be combined into
\begin{equation}
  p^\mu = i\partial^\mu ,
\end{equation}
raising the index on $\partial_\mu$ conveniently accounts for the minus
sign.  The plane wave $e^{-ik\cdot x}$ has momentum
\begin{equation}
  k^\mu ,
\end{equation}
since
\begin{equation}
  i\partial^\mu \bigl( e^{-ik\cdot x} \bigr) = k^\mu e^{-ik\cdot x} .
\end{equation}
The notation ``h.c.'' denotes the Hermitian conjugate.

Discussions of spin in quantum mechanics make use of the Pauli sigma
matrices:
\begin{equation}
  \sigma^1 = \begin{pmatrix}0&1\\1&0\end{pmatrix},\qquad
  \sigma^2 = \begin{pmatrix}0&-i\\i&0\end{pmatrix},\qquad
  \sigma^3 = \begin{pmatrix}1&0\\0&-1\end{pmatrix}.
\end{equation}
Products of these matrices satisfy the identity
\begin{equation}
  \sigma^i \sigma^j = \delta^{ij} + i\epsilon^{ijk} \sigma^k .
\end{equation}
It is convenient to define the linear combinations
$\sigma^{\pm} = \frac{1}{2}(\sigma^1 \pm i\sigma^2)$; then
\begin{equation}
  [\sigma^3, \sigma^\pm] = \pm 2 \sigma^\pm .
\end{equation}

\section*{Fourier Transforms and Distributions}

We will often make use of the Heaviside step function $\theta(x)$ and the
Dirac delta function $\delta(x)$, defined as follows:
\begin{equation}
  \theta(x) = \begin{cases}1 & x>0\\ 0 & x<0 \end{cases},\qquad
  \int_{-\infty}^{\infty} \!\! dx\ \delta(x) = 1 .
\end{equation}
The delta function in $n$ dimensions, denoted $\delta^{(n)}(\mathbf{x})$,
is zero everywhere except at $\mathbf{x}=0$ and satisfies
\begin{equation}
  \int d^n x\ \delta^{(n)}(\mathbf{x}) = 1 .
\end{equation}
In Fourier transforms the factors of $2\pi$ will always appear with the
momentum integral.  For example, in four dimensions,
\begin{equation}
  f(x) = \int \frac{d^4k}{(2\pi)^4}\; e^{-ik\cdot x} f(k),\qquad
  f(k) = \int d^4x\; e^{ik\cdot x} f(x) .
\end{equation}
(In three-dimensional transforms the signs in the exponents will be $+$
and $-$, respectively.)  The tilde on $f(k)$ will sometimes be omitted when
there is no potential for confusion.  The other important factors of $2\pi$
to remember appear in the identity
\begin{equation}
  \int d^4x\; e^{ik\cdot x} = (2\pi)^4 \delta^{(4)}(k) .
\end{equation}

\section*{Electrodynamics}

We use the Heaviside--Lorentz conventions, in which the factors of $4\pi$
appear in Coulomb's law and the fine-structure constant rather than in
Maxwell's equations.  Thus the Coulomb potential of a point charge $Q$ is
\begin{equation}
  \Phi = \frac{Q}{4\pi r} ,
\end{equation}
and the fine-structure constant is
\begin{equation}
  \alpha = \frac{e^2}{4\pi} = \frac{e^2}{4\pi \hbar c} = \frac{1}{137} .
\end{equation}
The symbol $e$ stands for the charge of the electron, a negative quantity
(although the sign rarely matters).  We generally work with the
relativistic form of Maxwell's equations:
\begin{equation}
  \partial_\mu F^{\mu\nu} = e J^\nu ,
\end{equation}
\begin{equation}
  \epsilon^{\mu\nu\rho\sigma}\ \partial_\nu F_{\rho\sigma} = 0 ,
\end{equation}
where
\begin{equation}
  A^\mu = (\Phi,\ \mathbf{A}),\qquad
  F^{\mu\nu} = \partial^\mu A^\nu - \partial^\nu A^\mu
\end{equation}
and we have extracted the $e$ from the 4-vector current density
\begin{equation}
  J^\mu = (\rho,\ \mathbf{J}) .
\end{equation}

\section*{Dirac Equation}

Some of our conventions differ from those of Bjorken and Drell (1964,
1965) and other texts: We use a \emph{chiral basis} for Dirac matrices,
and \emph{relativistic normalization} for Dirac spinors.  These
conventions are introduced in Sections~3.2 and~3.3, and are summarized in
the Appendix.

% ======================================================================
\chapter*{Editor's Foreword}
\addcontentsline{toc}{chapter}{Editor's Foreword}

The problem of communicating in a coherent fashion recent developments in
the most exciting and active fields of physics continues to be with us.
The enormous growth in the number of physicists has tended to make the
familiar channels of communication considerably less effective.  It has
become increasingly difficult for experts in a given field to keep up with
the current literature; the novice can only be confused.  What is needed
is both a consistent account of a field and the presentation of a definite
``point of view'' concerning it.  Formal monographs cannot meet such a need
in a rapidly developing field, while the review article seems to have
fallen into disfavor.  Indeed, it would seem that the people most actively
engaged in developing a given field are the people least likely to write
at length about it.

\textit{Frontiers in Physics} was conceived in 1961 in an effort to
improve the situation in several ways.  Leading physicists frequently give
a series of lectures, a graduate seminar, or a graduate course in their
special fields of interest.  Such lectures serve to summarize the present
status of a rapidly developing field and may well constitute the only
coherent account available at the time.  Often, notes on lectures exist
(prepared by the lecturer, by graduate students, or by postdoctoral
fellows) and are distributed in photocopied form on a limited basis.  One
of the principal purposes of the \textit{Frontiers in Physics} series is
to make such notes available to a wider audience of physicists.

As \textit{Frontiers in Physics} has evolved, a second category of book,
the informal text/monograph, an intermediate step between lecture notes
and formal texts or monographs, has played an increasingly important role
in the series.  In an informal text or monograph an author has reworked
his/her lecture notes to the point at which the manuscript represents a
coherent summation of a newly developed field, complete with references
and problems, suitable for either classroom teaching or individual study.

During the past two decades significant advances have been made in both
the conceptual framework of quantum field theory and its application to
condensed matter physics and elementary particle physics.  Given the fact
that the study of quantum field theory has become an essential part of the
education of graduate students in physics, a textbook which makes these
recent developments accessible to the novice, while not neglecting the
basic concepts, is highly desirable.  Michael Peskin and Daniel Schroeder
have written just such a book, describing in lucid fashion quantum
electrodynamics, renormalization, and non-Abelian gauge theories while
offering the reader a taste of what is to come.  It is therefore quite
appropriate to include this very polished text/monograph in the
\textit{Frontiers in Physics} series, and it gives me pleasure to welcome
them to the ranks of its authors.

\vspace{2em}
\noindent Aspen, Colorado\\
August 1995

\vspace{1em}
\noindent David Pines
